\section{Discrete formal smoothness and Cohen lifts (Part A)}\label{sec:A}

We give the discrete topology to every ring in Sections~\ref{subsec:A1}
and~\ref{subsec:A2}.  In particular, ordinary K\"ahler differentials are
not replaced by their completion.  A map of commutative rings $R\to S$ is
formally smooth if every $R$-algebra map $S\to C/J$, for an $R$-algebra
$C$ and an ideal $J$ satisfying $J^2=0$, lifts to an $R$-algebra map
$S\to C$.  The test algebra and its module $J$ are unrestricted.

\begin{lemma}[Proved: two lifting consequences]\label{SA23-A-01}
If $S$ is formally smooth over $R$, then $\Omega_{S/R}$ is projective
over $S$.  An $R$-algebra retract of a formally smooth $R$-algebra is
formally smooth over $R$.
\end{lemma}
\begin{proof}
For a surjective homomorphism $M\to N$ of $S$-modules, the map of
square-zero $R$-algebras $S\oplus M\to S\oplus N$ is surjective.
An $R$-derivation $d:S\to N$ gives the $R$-algebra map
$s\mapsto(s,d(s))$.  A lift has first coordinate $s$, so its second
coordinate is a derivation $S\to M$ lifting $d$.  By the defining
universal property of $\Omega_{S/R}$, its Hom functor preserves
surjections, which is precisely projectivity.  This is also the
consequence of the split conormal criterion in
\cite[Tag 00TH, Lemma 10.138.7]{rootStacksSmooth}.
For the second assertion, write $T\xrightarrow{i}S\xrightarrow{r}T$
with $ri=1_T$.  A test map $T\to C/J$ composed with $r$ lifts on $S$;
restriction along $i$ supplies the required lift on $T$.
\end{proof}

\begin{lemma}[Proved: a separation obstruction]\label{SA23-A-02}
Let $R$ be a ring, $t\in R$, and assume $\bigcap_{N\geq0}t^NR=0$.
Every projective $R$-module $P$ satisfies
$\bigcap_{N\geq0}t^NP=0$.
\end{lemma}
\begin{proof}
Choose a split injection $P\to R^{(J)}$ into a free module.  The image
of an element in every $t^NP$ has each coordinate in every $t^NR$,
so all its coordinates vanish.  The injection then kills only zero.
No finite-generation assumption on $P$ is used.
\end{proof}

\subsection{The characteristic-zero answer}\label{subsec:A1}

\begin{lemma}[Proved: the formal exponential is transcendental]\label{SA23-A-03}
For every field $K$ of characteristic zero, the element
$f=\sum_{m\geq0}t^m/m!$ of $K[[t]]$ is transcendental over $K(t)$.
\end{lemma}
\begin{proof}
Suppose it is algebraic, and let
$P(Y)=Y^d+a_{d-1}(t)Y^{d-1}+\cdots+a_0(t)$ in $K(t)[Y]$
be its monic minimal polynomial.  Here $d>0$ and $a_0\ne0$, since
$f\ne0$.  The usual coefficientwise derivative on $K((t))$ is a
$K$-derivation and $f'=f$.  Differentiating $P(f)=0$ gives
$Q(f)=0$, where
\[
 Q(Y)=dY^d+\sum_{j=0}^{d-1}\bigl(a_j'(t)+j a_j(t)\bigr)Y^j.
\]
The polynomial $Q-dP$ has degree less than $d$ and vanishes at $f$;
minimality implies that it is zero.  Its constant coefficient gives
$a_0'=d a_0$.  No nonzero rational function has this property: if
$a_0=u/v$ with nonzero polynomials $u,v\in K[t]$, then
$u'v-uv'=d uv$, whose left side has degree at most
$\deg u+\deg v-1$ and whose right side has degree
$\deg u+\deg v$.  This contradiction proves the assertion.
\end{proof}

\begin{theorem}[Proved: A.1, with no restriction on $\Omega_{K/k}$]\label{SA23-A-04}
Let $k\subset K$ be fields of characteristic zero, let $n\geq1$, and
put $A=K[[t_1,\ldots,t_n]]$.  The ordinary $A$-module $\Omega_{A/k}$
is not projective, and therefore is not free.  Consequently $A$ is
not formally smooth over $k$ for the discrete topology.  In
particular this holds for $\mathbf Q[[t]]/\mathbf Q$ and for every
case of A.1 with finite-dimensional $\Omega_{K/k}$.
\end{theorem}
\begin{proof}
Let $t=t_1$ and $f=\exp(t)\in A$.  Write $d:A\to\Omega_{A/k}$
for the universal derivation and put $\eta=df-f\,dt$.
For each integer $N\geq1$ write
\[
 f=P_N+t^Ng_N,\qquad
 P_N=\sum_{m=0}^{N-1}\frac{t^m}{m!},\qquad g_N\in K[[t]].
\]
All coefficients of $P_N$ belong to the prime field and hence to $k$.
Formal differentiation, which is an operation on power series and is
not an assertion about arbitrary derivations, gives
$f=P_N'+Nt^{N-1}g_N+t^Ng_N'$.  The finite product rule for $d$ now
gives the exact identity
\[
 \eta=t^N\bigl(dg_N-g_N'\,dt\bigr).
\]
Thus $\eta$ belongs to every $t^N\Omega_{A/k}$.

We show that $\eta\ne0$.  The elements $t_2,\ldots,t_n$ are
algebraically independent over $K((t))$: a polynomial relation, viewed
as a series in these variables with coefficients in $K((t))$, has
all coefficients zero.  Lemma~\ref{SA23-A-03} therefore shows that
$t_1,\ldots,t_n,f$ are algebraically independent over $K$ in
$F=\operatorname{Frac}(A)$.  Extend them to a transcendence basis
$\mathcal T$ of $F/K$.  On the rational function field $K(\mathcal T)$
there is the $K$-derivation $D$ taking $f$ to $1$ and all other basis
elements to $0$.  Since $F/K(\mathcal T)$ is algebraic and
separable in characteristic zero, $D$ extends to $F$.  Explicitly,
for an algebraic element $z$ with minimal polynomial $q$, its forced
value is $D(z)=-D(q)(z)/q'(z)$; this constructs the unique extension
on each finite subextension, and uniqueness makes the extensions
agree on overlaps.  Restriction gives a $k$-derivation $A\to F$.
The induced $A$-linear map $\Omega_{A/k}\to F$ sends $\eta$ to $1$.
Hence $\eta\ne0$.  Since $A$ is $t$-adically separated,
Lemma~\ref{SA23-A-02} rules out projectivity.  Apply
Lemma~\ref{SA23-A-01} to rule out formal smoothness.
\end{proof}

\begin{remark}[Refuted reading and proved scope]\label{SA23-A-05}
This proves the positive-dimensional presumption of cote 22/023
without assuming finite transcendence degree of $K/k$.  The two
supplied witnesses give either $\dim A>0$ or a doubtful
$\dim A\geq0$.  If the latter includes $n=0$, that reading is false:
every characteristic-zero field extension is formally smooth and
its differentials form a free vector space
\cite[Tag 031U, Proposition 10.158.9]{rootStacksFields}.
The same-field non-smoothness assertion is located by secondary
literature at Matsumura, \emph{Commutative Ring Theory}, p.~215;
that page has not yet been opened, so we rely on the complete proof
above and make no novelty claim.
\end{remark}

\subsection{The complete characteristic-\texorpdfstring{$p$}{p} criterion}\label{subsec:A2}

\begin{lemma}[Proved: infinitely many coefficient differentials]\label{SA23-A-06}
Let $k\subset K$ be fields of any characteristic and assume
$\dim_K\Omega_{K/k}$ is infinite.  For every $n\geq1$ the module
$\Omega_{K[[t_1,\ldots,t_n]]/k}$ is not projective.
\end{lemma}
\begin{proof}
Put $A=K[[t_1,\ldots,t_n]]$, $\mathfrak m=(t_1,\ldots,t_n)$, and
$t=t_1$.  If the module were projective, choose a split injection
$j:\Omega_{A/k}\to M=A^{(J)}$.  The augmentation $A\to K$ and the
constant inclusion $K\to A$ imply that the map
$\Omega_{K/k}\to\Omega_{A/k}\otimes_AK$ is split injective: the
reverse map is induced by the augmentation, and their composite
is the identity on each $da$.  Since $j$ remains injective after
tensoring with $K$, the vectors $j(da)\bmod\mathfrak mM$, for
$a\in K$, span an infinite-dimensional subspace of $M/\mathfrak mM$.

For $q\in J$, let $D_q:A\to A$ be the coordinate of $j\circ d$.
The vector $j(dt)$ has finite support.  Inductively choose
$a_i\in K$ and distinct $q_i\in J$, for $i\geq1$, such that $q_i$
is outside the supports of $j(dt),j(da_1),\ldots,j(da_{i-1})$ and
$D_{q_i}(a_i)\notin\mathfrak m$.  The infinite-dimensionality just
proved permits every step.  Each $D_{q_i}$ kills $t$, and hence
preserves $t^NA$ for all $N$ by the finite product rule.
For $f=\sum_{i\geq1}a_it^i\in A$, the identities modulo $t^NA$
therefore give
\[
 D_{q_i}(f)=\sum_{l\geq i} D_{q_i}(a_l)t^l.
\]
This equality follows from $t$-adic separation of the target $A$;
we have not assumed continuity of derivations into arbitrary
modules.  Its coefficient at order $i$ modulo $t$ is nonzero,
because $D_{q_i}(a_i)$ is a unit.  Thus $j(df)$ has infinitely many
nonzero coordinates, contrary to $j(df)\in A^{(J)}$.
\end{proof}

\begin{lemma}[Proved: the Frobenius coefficient obstruction]\label{SA23-A-07}
Let $k\subset K$ have characteristic $p>0$, and suppose
$[kK^p:K^p]$ is infinite.  Then $\Omega_{K[[t]]/k}$ contains a
nonzero element in $\bigcap_{N\geq0}t^N\Omega_{K[[t]]/k}$.
\end{lemma}
\begin{proof}
The $K^p$-linear span of $k$ inside $K$ is the field $kK^p$.
Indeed the span is a subring, and a nonzero element $a$ in it has
$a^p\in K^p$, so $a^{-1}=a^{p-1}/a^p$ belongs to the span.
We can therefore choose $a_i\in k$, $i\geq1$, linearly independent
over $K^p$.  Put $f=\sum_{i\geq1}a_it^{pi}$ and $\eta=df$.
Subtracting a finite initial sum and applying $d$ gives
$df=t^{pN}dg_N$ for a tail $g_N\in K[[t]]$; the coefficient
differentials and $d(t^{pN})$ are zero.  Thus $\eta$ lies in every
$t^N\Omega$.

Let $F=K((t))$.  For a characteristic-$p$ field extension $F/k$,
\[
 \ker(d:F\to\Omega_{F/k})=kF^p.
\]
Here is the elementary justification.  Set $E=kF^p$; every
$k$-derivation kills $E$.  Choose a $p$-basis $(b_j)$ of $F/E$ by
maximally enlarging a family with independent monomials
$\prod b_j^{e_j}$, $0\leq e_j<p$.  Maximality makes these monomials
span $F$, since adjoining an element outside the span would enlarge
an exponent-one purely inseparable extension by degree $p$.
Each element has a unique finite expansion in these monomials.
The formal partial derivatives with respect to $b_j$ are
$E$-derivations, because $b_j^p\in E$.  An element killed by all
of them has only the constant monomial in its expansion, and lies
in $E$.  This proves the displayed formula.

We have $F^p=K^p((t^p))$.  Moreover
\[
 KF^p=\bigcup_{E_0} E_0((t^p)),
\]
where $E_0$ runs through the intermediate fields
$K^p\subset E_0\subset K$ finite over $K^p$.  To see the equality,
each element of the compositum uses finitely many constants from
$K$; they generate such an $E_0$ since their $p$th powers belong to
$K^p$.  Conversely a finite $K^p$-basis of $E_0$ expands every
Laurent series with $E_0$-coefficients into a finite sum of that
basis times series in $K^p((t^p))$.
The coefficients of $f$ cannot all lie in any such finite-dimensional
$E_0$, so $f\notin KF^p$, and hence $f\notin kF^p$.
Thus $df$ has nonzero image in $\Omega_{F/k}$ and is itself nonzero.
\end{proof}

\begin{theorem}[Proved: A.2, exact criterion]\label{SA23-A-08}
Let $k\subset K$ be fields of characteristic $p>0$ and let $n\geq1$.
The following conditions are equivalent:
\begin{enumerate}
\item $K[[t_1,\ldots,t_n]]$ is formally smooth over $k$ for the
discrete topology.
\item $K/k$ is separable and $[K:K^p]<\infty$.
\end{enumerate}
If (2) holds and a $p$-basis $c_1,\ldots,c_r$ of $k$ extends to
a $p$-basis $c_1,\ldots,c_r,b_1,\ldots,b_s$ of $K$, then
$\Omega_{K[[t_1,\ldots,t_n]]/k}$ is free on
$db_1,\ldots,db_s,dt_1,\ldots,dt_n$.
\end{theorem}
\begin{proof}
Write $A=K[[t_1,\ldots,t_n]]$.  If $A/k$ is formally smooth, the
coefficient field $K$ is a $k$-algebra retract, so $K/k$ is formally
smooth.  This is equivalent to separability by
\cite[Tag 031U, Proposition 10.158.9]{rootStacksFields}.
Suppose that $[K:K^p]$ is infinite.  If $\dim_K\Omega_{K/k}$ is
infinite, Lemmas~\ref{SA23-A-01} and~\ref{SA23-A-06} give a
contradiction.  Otherwise a finite $p$-basis of $K/kK^p$ shows that
$[K:kK^p]$ is finite, so $[kK^p:K^p]$ must be infinite.  Lemmas
\ref{SA23-A-02} and~\ref{SA23-A-07} rule out formal smoothness
of $K[[t_1]]/k$.  This one-variable algebra is a retract of $A$
by setting $t_2,\ldots,t_n$ to zero, so it also rules out formal
smoothness of $A/k$.

Conversely suppose (2).  Separability means that every
$p$-independent family in $k$ remains $p$-independent over $K^p$;
equivalently the field map
$k\otimes_{k^p}K^p\to K$ is injective.  Thus a $p$-basis of $k$
is finite and extends to a finite $p$-basis $c_1,\ldots,c_r,b_1,
\ldots,b_s$ of $K$.  Separating the exponents of the variables
modulo $p$, and expanding each coefficient in this finite basis,
shows that $c,b,t$ is a $p$-basis of $A$ over $A^p$.
The finiteness is essential: only finitely many Frobenius-monomial
components have to be assembled into power series.
Consequently the natural map $k\otimes_{k^p}A^p\to A$ is
injective, with image $R=kA^p$, and there is an isomorphism
\[
 A\simeq
 R[X_1,\ldots,X_{s+n}]/(X_i^p-u_i^p\mid1\leq i\leq s+n),
 \qquad (u_1,\ldots,u_{s+n})=(b_1,\ldots,b_s,t_1,\ldots,t_n).
\]
Its injectivity and surjectivity both follow from the same unique
monomial expansion.

Now let $C$ be a $k$-algebra, $J\subset C$ with $J^2=0$, and
$h:A\to C/J$ a $k$-algebra map.  There is a canonical ring map
$v:A^p\to C$: if $a^p\in A^p$, choose any lift $\widetilde h(a)$
of $h(a)$ and put $v(a^p)=\widetilde h(a)^p$.
Frobenius on $A$ is injective, so $a$ is unique; two choices of
lift differ by an element of $J$, whose $p$th power is zero.
The characteristic-$p$ binomial identity proves that $v$ is additive
and multiplicative.  On $k^p$ it agrees with the given structure
map to $C$.  It therefore extends to $R=k\otimes_{k^p}A^p$ as a
$k$-algebra map.  Choose arbitrary lifts $\widetilde h(u_i)$.
Their $p$th powers are $v(u_i^p)$, so the presentation above gives
a $k$-algebra map $A\to C$ lifting $h$.

Finally every derivation of $A/k$ kills $R$, and the presentation
has zero differentials for all its defining relations.  The
universal property therefore identifies $\Omega_{A/k}$ with the
free $A$-module on the displayed $du_i$.
\end{proof}

\begin{corollary}[Refuted: finite relative $p$-basis without a perfect base]\label{SA23-A-09}
The assertion on cote 22/022 that a finite $p$-basis of $K/k$ always
extends by the $t_i$ to a $p$-basis of $K[[t_1,\ldots,t_n]]/k$
is false if arbitrary imperfect $k$ is allowed.  Take
$k=K=\mathbf F_p(s_1,s_2,\ldots)$, where the $s_i$ are algebraically
independent.  The relative $p$-basis is empty, but for $n\geq1$
the power series ring is not formally smooth over $k$.
If $k$ is perfect and $K/k$ is separable, the manuscript's
finite-relative-$p$-basis criterion is correct, and the infinite
case has a negative answer.
\end{corollary}
\begin{proof}
For the specified rational function field the monomials in finitely
many $s_i$ with exponents less than $p$ form a basis over $K^p$;
in particular $[K:K^p]$ is infinite.  Apply
Theorem~\ref{SA23-A-08}.  More concretely,
$d(\sum_{i\geq1}s_it^{pi})$ is the nonzero infinitely divisible
element in Lemma~\ref{SA23-A-07}.  If $k$ is perfect, then
$k\subset K^p$, so finite relative $p$-rank is exactly finite
absolute $p$-rank.  This treats both the perfect-base branch and
the broader separable-residue-extension branch of the doubtful
introductory wording on scan 022.
\end{proof}
